\documentclass[reqno]{amsart}
\usepackage[T1]{fontenc}
\usepackage{amsmath,amssymb,amsthm}
\usepackage{xr-hyper}
\usepackage[hidelinks]{hyperref}
\externaldocument[dt-]{20261125T002000--definite-theta-lifts__workbook}
\externaldocument[ps-]{20261125T010500--poincare-series-and-near-orthogonality__workbook}
\externaldocument[as-]{20261125T021500--assembly-at-one-scale__workbook}

% Chapter number; \ref{chap} gives it.
\newcounter{chap}
% Chapter numbers, generated from 20261125T012000--grh-workbook__guide.org
% by scripts/make-chapters.py.  Do not edit: change the order of the
% index and rerun the script (./preview-site.sh and CI do this).
%
% Usage in a sheet, after \newcounter{chap}:
%   % Chapter numbers, generated from 20261125T012000--grh-workbook__guide.org
% by scripts/make-chapters.py.  Do not edit: change the order of the
% index and rerun the script (./preview-site.sh and CI do this).
%
% Usage in a sheet, after \newcounter{chap}:
%   % Chapter numbers, generated from 20261125T012000--grh-workbook__guide.org
% by scripts/make-chapters.py.  Do not edit: change the order of the
% index and rerun the script (./preview-site.sh and CI do this).
%
% Usage in a sheet, after \newcounter{chap}:
%   \input{chapters}\chapnumber{<sheet ID>}
\makeatletter
\providecommand\chapters@set[2]{\@namedef{chapters@#1}{#2}}
\chapters@set{20261125T001000}{1}% zeros-to-ramp-sums
\chapters@set{20261125T001100}{2}% inert-primes-and-estimate-b
\chapters@set{20261125T005000}{3}% galois-fourier-and-averaging
\chapters@set{20261125T005500}{4}% columns-from-a-sphere
\chapters@set{20261125T000000}{5}% slope-method
\chapters@set{20261126T030100}{6}% a-twist-for-each-prime
\chapters@set{20261125T003000}{7}% the-chain-and-the-square-root
\chapters@set{20261125T010000}{8}% rank-lemma-and-zero-estimate
\chapters@set{20261125T010500}{9}% poincare-series-and-near-orthogonality
\chapters@set{20261124T223000}{10}% elliptic-curves-for-the-proof
\chapters@set{20261125T009500}{11}% fock-model-and-theta-intertwiners
\chapters@set{20261124T223100}{12}% complex-multiplication-and-inert-reduction
\chapters@set{20261125T008000}{13}% one-cm-determinant
\chapters@set{20261124T214500}{14}% supersingular-theta-divisibility
\chapters@set{20261124T214600}{15}% the-isotropic-plane
\chapters@set{20261125T013700}{16}% pairings-on-a-formal-group-of-height-two
\chapters@set{20261125T013500}{17}% hecke-nondegeneracy-on-the-tree
\chapters@set{20261126T030200}{18}% the-bound-is-attained
\chapters@set{20261125T009000}{19}% kronecker-limit-and-chowla-selberg
\chapters@set{20261125T002000}{20}% definite-theta-lifts
\chapters@set{20261125T015500}{21}% siegel-weil-in-the-smallest-case
\chapters@set{20261125T002100}{22}% definite-theta-lifts-and-norms
\chapters@set{20261125T016500}{23}% hecke-operators-and-theta-lifts
\chapters@set{20261125T016700}{24}% divisibility-along-the-chain
\chapters@set{20261125T004000}{25}% denominators-and-sha
\chapters@set{20261125T004100}{26}% tate-twists-and-sha
\chapters@set{20261126T020430}{27}% the-minus-part-by-gauss-sums
\chapters@set{20261125T014100}{28}% congruence-modules
\chapters@set{20261125T014500}{29}% eisenstein-congruences-and-selmer-classes
\chapters@set{20261125T014700}{30}% pseudocharacters-and-the-length-bound
\chapters@set{20261125T011000}{31}% galois-representations-black-box
\chapters@set{20261126T030300}{32}% the-data
\chapters@set{20261126T020210}{33}% bessel-poincare-series-and-their-gram-operators
\chapters@set{20261126T030400}{34}% the-seesaw-identity
\chapters@set{20261126T030500}{35}% divisibility-of-the-composite-lift
\chapters@set{20261126T030600}{36}% theorem-a-short-forms
\chapters@set{20261126T030700}{37}% the-determinant-of-a-chain
\chapters@set{20261126T030800}{38}% the-cost-of-one-lift
\chapters@set{20261126T030900}{39}% denominators-of-projections-to-lifts
\chapters@set{20261126T031000}{40}% twists-with-few-congruences
\chapters@set{20261126T031100}{41}% theorem-c-and-the-main-inequality
\providecommand\chapnumber[1]{%
  \@ifundefined{chapters@#1}%
    {\PackageError{chapters}{Sheet #1 is not listed in the index}%
      {List it in 20261125T012000--grh-workbook__guide.org\MessageBreak
       and run scripts/make-chapters.py.}}%
    {\edef\thechap{\@nameuse{chapters@#1}}}}
\makeatother
\chapnumber{<sheet ID>}
\makeatletter
\providecommand\chapters@set[2]{\@namedef{chapters@#1}{#2}}
\chapters@set{20261125T001000}{1}% zeros-to-ramp-sums
\chapters@set{20261125T001100}{2}% inert-primes-and-estimate-b
\chapters@set{20261125T005000}{3}% galois-fourier-and-averaging
\chapters@set{20261125T005500}{4}% columns-from-a-sphere
\chapters@set{20261125T000000}{5}% slope-method
\chapters@set{20261126T030100}{6}% a-twist-for-each-prime
\chapters@set{20261125T003000}{7}% the-chain-and-the-square-root
\chapters@set{20261125T010000}{8}% rank-lemma-and-zero-estimate
\chapters@set{20261125T010500}{9}% poincare-series-and-near-orthogonality
\chapters@set{20261124T223000}{10}% elliptic-curves-for-the-proof
\chapters@set{20261125T009500}{11}% fock-model-and-theta-intertwiners
\chapters@set{20261124T223100}{12}% complex-multiplication-and-inert-reduction
\chapters@set{20261125T008000}{13}% one-cm-determinant
\chapters@set{20261124T214500}{14}% supersingular-theta-divisibility
\chapters@set{20261124T214600}{15}% the-isotropic-plane
\chapters@set{20261125T013700}{16}% pairings-on-a-formal-group-of-height-two
\chapters@set{20261125T013500}{17}% hecke-nondegeneracy-on-the-tree
\chapters@set{20261126T030200}{18}% the-bound-is-attained
\chapters@set{20261125T009000}{19}% kronecker-limit-and-chowla-selberg
\chapters@set{20261125T002000}{20}% definite-theta-lifts
\chapters@set{20261125T015500}{21}% siegel-weil-in-the-smallest-case
\chapters@set{20261125T002100}{22}% definite-theta-lifts-and-norms
\chapters@set{20261125T016500}{23}% hecke-operators-and-theta-lifts
\chapters@set{20261125T016700}{24}% divisibility-along-the-chain
\chapters@set{20261125T004000}{25}% denominators-and-sha
\chapters@set{20261125T004100}{26}% tate-twists-and-sha
\chapters@set{20261126T020430}{27}% the-minus-part-by-gauss-sums
\chapters@set{20261125T014100}{28}% congruence-modules
\chapters@set{20261125T014500}{29}% eisenstein-congruences-and-selmer-classes
\chapters@set{20261125T014700}{30}% pseudocharacters-and-the-length-bound
\chapters@set{20261125T011000}{31}% galois-representations-black-box
\chapters@set{20261126T030300}{32}% the-data
\chapters@set{20261126T020210}{33}% bessel-poincare-series-and-their-gram-operators
\chapters@set{20261126T030400}{34}% the-seesaw-identity
\chapters@set{20261126T030500}{35}% divisibility-of-the-composite-lift
\chapters@set{20261126T030600}{36}% theorem-a-short-forms
\chapters@set{20261126T030700}{37}% the-determinant-of-a-chain
\chapters@set{20261126T030800}{38}% the-cost-of-one-lift
\chapters@set{20261126T030900}{39}% denominators-of-projections-to-lifts
\chapters@set{20261126T031000}{40}% twists-with-few-congruences
\chapters@set{20261126T031100}{41}% theorem-c-and-the-main-inequality
\providecommand\chapnumber[1]{%
  \@ifundefined{chapters@#1}%
    {\PackageError{chapters}{Sheet #1 is not listed in the index}%
      {List it in 20261125T012000--grh-workbook__guide.org\MessageBreak
       and run scripts/make-chapters.py.}}%
    {\edef\thechap{\@nameuse{chapters@#1}}}}
\makeatother
\chapnumber{<sheet ID>}
\makeatletter
\providecommand\chapters@set[2]{\@namedef{chapters@#1}{#2}}
\chapters@set{20261125T001000}{1}% zeros-to-ramp-sums
\chapters@set{20261125T001100}{2}% inert-primes-and-estimate-b
\chapters@set{20261125T005000}{3}% galois-fourier-and-averaging
\chapters@set{20261125T005500}{4}% columns-from-a-sphere
\chapters@set{20261125T000000}{5}% slope-method
\chapters@set{20261126T030100}{6}% a-twist-for-each-prime
\chapters@set{20261125T003000}{7}% the-chain-and-the-square-root
\chapters@set{20261125T010000}{8}% rank-lemma-and-zero-estimate
\chapters@set{20261125T010500}{9}% poincare-series-and-near-orthogonality
\chapters@set{20261124T223000}{10}% elliptic-curves-for-the-proof
\chapters@set{20261125T009500}{11}% fock-model-and-theta-intertwiners
\chapters@set{20261124T223100}{12}% complex-multiplication-and-inert-reduction
\chapters@set{20261125T008000}{13}% one-cm-determinant
\chapters@set{20261124T214500}{14}% supersingular-theta-divisibility
\chapters@set{20261124T214600}{15}% the-isotropic-plane
\chapters@set{20261125T013700}{16}% pairings-on-a-formal-group-of-height-two
\chapters@set{20261125T013500}{17}% hecke-nondegeneracy-on-the-tree
\chapters@set{20261126T030200}{18}% the-bound-is-attained
\chapters@set{20261125T009000}{19}% kronecker-limit-and-chowla-selberg
\chapters@set{20261125T002000}{20}% definite-theta-lifts
\chapters@set{20261125T015500}{21}% siegel-weil-in-the-smallest-case
\chapters@set{20261125T002100}{22}% definite-theta-lifts-and-norms
\chapters@set{20261125T016500}{23}% hecke-operators-and-theta-lifts
\chapters@set{20261125T016700}{24}% divisibility-along-the-chain
\chapters@set{20261125T004000}{25}% denominators-and-sha
\chapters@set{20261125T004100}{26}% tate-twists-and-sha
\chapters@set{20261126T020430}{27}% the-minus-part-by-gauss-sums
\chapters@set{20261125T014100}{28}% congruence-modules
\chapters@set{20261125T014500}{29}% eisenstein-congruences-and-selmer-classes
\chapters@set{20261125T014700}{30}% pseudocharacters-and-the-length-bound
\chapters@set{20261125T011000}{31}% galois-representations-black-box
\chapters@set{20261126T030300}{32}% the-data
\chapters@set{20261126T020210}{33}% bessel-poincare-series-and-their-gram-operators
\chapters@set{20261126T030400}{34}% the-seesaw-identity
\chapters@set{20261126T030500}{35}% divisibility-of-the-composite-lift
\chapters@set{20261126T030600}{36}% theorem-a-short-forms
\chapters@set{20261126T030700}{37}% the-determinant-of-a-chain
\chapters@set{20261126T030800}{38}% the-cost-of-one-lift
\chapters@set{20261126T030900}{39}% denominators-of-projections-to-lifts
\chapters@set{20261126T031000}{40}% twists-with-few-congruences
\chapters@set{20261126T031100}{41}% theorem-c-and-the-main-inequality
\providecommand\chapnumber[1]{%
  \@ifundefined{chapters@#1}%
    {\PackageError{chapters}{Sheet #1 is not listed in the index}%
      {List it in 20261125T012000--grh-workbook__guide.org\MessageBreak
       and run scripts/make-chapters.py.}}%
    {\edef\thechap{\@nameuse{chapters@#1}}}}
\makeatother
\chapnumber{20261125T015500}
\renewcommand{\thesection}{\thechap.\arabic{section}}

\theoremstyle{definition}
\newtheorem{problem}{Problem}[section]
\theoremstyle{plain}
\newtheorem{theorem}{Theorem}[section]

\newcommand{\Z}{\mathbb{Z}}
\newcommand{\Q}{\mathbb{Q}}
\newcommand{\R}{\mathbb{R}}
\newcommand{\C}{\mathbb{C}}
\newcommand{\F}{\mathbb{F}}
\newcommand{\ii}{\mathrm{i}}
\newcommand{\Oc}{\mathcal{O}}
\newcommand{\Mf}{\mathcal{M}}
\newcommand{\Ac}{\mathcal{A}}
\newcommand{\Cl}{\mathrm{Cl}}
\newcommand{\Sym}{\operatorname{Sym}}
\newcommand{\Gram}{\operatorname{Gram}}

\setlength{\emergencystretch}{2em}

\title{Siegel--Weil in the smallest case}
\author{}
\date{November 25, 2026}

\begin{document}
\maketitle
\refstepcounter{chap}\label{chap}

\paragraph{Prerequisites.}
Modular forms for \(\mathrm{SL}_2(\Z)\), theta series of even unimodular lattices, and the mass formula, as in \cite[Ch.~V, \S2.3 and Ch.~VII, \S\S2--6]{SeA}.  Quadratic forms over finite fields \cite[Ch.~IV, \S1.7]{SeA} and Hensel's lemma \cite[Ch.~II, \S2.2]{SeA}.  For \S\ref{sec:hermitian}: the space of modular forms of weight \(n\) on \(\Gamma_0(4)\) with character \(\chi^n\), its dimension and its basis in terms of \(\theta\) \cite[Ch.~IV, \S1]{Ko}, and its Eisenstein series \cite[\S4.5]{DS}.  From Chapter~\ref{dt-chap}, \emph{Masses and the count of forms}: the class sets \(\Cl_n\), the groups \(\Gamma_L\), the masses \(M_n\), Problems~\ref{dt-pr:aut}, \ref{dt-pr:class-one} and~\ref{dt-pr:rank-five}, and the forms of \S\ref{dt-sec:forms} with Problem~\ref{dt-pr:sym-norm}.

A class set is a finite set of lattices, and an individual lattice in it can be complicated.  The average over the class set, each lattice weighted by the inverse of the order of its automorphism group, is simple: this is Siegel's theorem, and in Weil's form it says that the mass-weighted average of theta series is an Eisenstein series.  We check it where it can be checked by hand, then on the class set of Chapter~\ref{dt-chap} in rank \(5\), and then say why the Gram matrices of the theta lifts in the argument are controlled by it.


\section{Even unimodular lattices}\label{sec:even}

Let \(L\) be an even unimodular positive definite lattice of rank \(m\), with inner product \(x \cdot x\), and let \(r_L(j)\) be the number of \(x \in L\) with \(x \cdot x = 2j\).  Then \(8 \mid m\), and \(\theta_L(\tau) := \sum_{j \ge 0} r_L(j) q^j\), \(q = e^{2\pi\ii\tau}\), is a modular form of weight \(k := m/2\) for \(\mathrm{SL}_2(\Z)\) \cite[Ch.~VII, Theorem~8]{SeA}.  Let \(\Mf_k\) be the space of such forms and
\[
  E_k := 1 - \frac{2k}{B_k} \sum_{j \ge 1} \sigma_{k-1}(j) q^j, \qquad \Delta := q\prod_{j \ge 1}(1-q^j)^{24} = \sum_{j \ge 1} \tau(j) q^j ,
\]
so that \(E_4 = 1 + 240\sum \sigma_3(j)q^j\), \(E_8 = 1 + 480\sum\sigma_7(j)q^j\), \(E_{12} = 1 + \frac{65520}{691}\sum\sigma_{11}(j)q^j\).  (Serre writes \(E_2, E_4, E_6\) for these, indexing by half the weight.)  The spaces \(\Mf_4\) and \(\Mf_8\) are spanned by \(E_4\) and \(E_8\), and \(\Mf_{12}\) by \(E_{12}\) and \(\Delta\) \cite[Ch.~VII, \S3.2]{SeA}.  We write \(E_8\) also for the root lattice; the context decides.  Let \(D_{16}^+ := D_{16} \cup (D_{16} + (\tfrac12, \dots, \tfrac12))\), where \(D_{16} = \{x \in \Z^{16} : \sum x_i \text{ even}\}\) \cite[Ch.~4, \S7]{CoS}.

\begin{problem}\label{pr:E8}
\(r_{E_8}(j) = 240\,\sigma_3(j)\) for \(j \ge 1\).  In particular \(E_8\) has \(240\) roots and \(2160\) vectors with \(x \cdot x = 4\).
\end{problem}

\begin{problem}\label{pr:rank16}
\(\theta_{E_8 \oplus E_8} = \theta_{D_{16}^+} = E_8\).  Both lattices have \(480\) roots and \(61920\) vectors with \(x\cdot x = 4\), and \(\sigma_7(j) = \sigma_3(j) + 120\sum_{i=1}^{j-1}\sigma_3(i)\sigma_3(j-i)\).
\end{problem}

\begin{problem}\label{pr:not-iso}
\(E_8 \oplus E_8\) and \(D_{16}^+\) are not isometric.
\end{problem}

So the theta series does not determine the lattice.  Milnor observed that the tori \(\R^{16}/(E_8\oplus E_8)\) and \(\R^{16}/D_{16}^+\) therefore have the same Laplace spectrum without being isometric \cite{Mi}.

\begin{problem}\label{pr:leech}
If \(\Lambda\) is an even unimodular lattice of rank \(24\) without roots (the Leech lattice \cite[Ch.~4, \S11]{CoS}), then \(\theta_\Lambda = E_{12} - \frac{65520}{691}\Delta\), and \(\Lambda\) has \(196560\) vectors with \(x\cdot x = 4\).
\end{problem}

In rank \(24\) the theta series is \(E_{12} + c_L\Delta\) with \(c_L = r_L(1) - \frac{65520}{691}\) \cite[Ch.~VII, \S6.6]{SeA}, and \(c_L\) varies with \(L\).  Let \(C_m\) be the set of isometry classes of even unimodular lattices of rank \(m\), let \(g_L := \lvert O(L) \rvert\), and let \(\mu_m := \sum_{L \in C_m} g_L^{-1}\) be the mass.  Siegel proved \cite{Si} that
\begin{equation}\label{eq:siegel}
  \sum_{L \in C_m} \frac{\theta_L}{g_L} = \mu_m\, E_{m/2},
\end{equation}
see \cite[Ch.~VII, \S6.5, Remark]{SeA}, and the Minkowski--Siegel mass formula \cite[Ch.~V, \S2.3]{SeA}
\[
  \mu_m = \frac{\lvert B_{k}\rvert}{2k} \prod_{j=1}^{k-1} \frac{\lvert B_{2j}\rvert}{4j}, \qquad m = 2k .
\]
One has \(C_8 = \{E_8\}\) with \(g_{E_8} = \lvert W(E_8)\rvert = 696729600\), and \(C_{16} = \{E_8 \oplus E_8, D_{16}^+\}\) with \(g = 2\,\lvert W(E_8)\rvert^2\) and \(g = 2^{15}\,16!\) respectively \cite[Ch.~V, \S2.3]{SeA}.

\begin{problem}\label{pr:mass}
The mass formula holds for \(m = 8\) and \(m = 16\).  In the average \eqref{eq:siegel} for \(m = 16\), \(E_8 \oplus E_8\) has weight \(286/691\) and \(D_{16}^+\) has weight \(405/691\).
\end{problem}

Problems~\ref{pr:E8} and~\ref{pr:rank16} prove \eqref{eq:siegel} for \(m = 8\) and \(16\): there every theta series is the Eisenstein series, because \(\dim \Mf_{m/2} = 1\).  In rank \(24\) the statement has content.

\begin{problem}\label{pr:niemeier}
For \(m = 24\), \eqref{eq:siegel} is equivalent to \(\sum_{L \in C_{24}} c_L/g_L = 0\).  Hence the mass-weighted mean number of roots of an even unimodular lattice of rank \(24\) is \(65520/691\); in particular \(\mu_{24}^{-1}\sum_L \theta_L/g_L\) is not the theta series of a lattice.
\end{problem}


\section{Local densities}\label{sec:densities}

Why should an average of theta series be an Eisenstein series?  Its coefficients are averages of representation numbers, and Siegel's theorem says that these are products of local quantities.  Let \(Q(x) := x\cdot x/2\) on \(L\).  For a prime \(p\) and \(j \ge 1\) the \emph{local density} is
\[
  \alpha_p(j) := \lim_{r \to \infty} p^{r(1-m)}\, \#\{x \in L/p^rL : Q(x) \equiv j \pmod{p^r}\},
\]
and the \emph{archimedean density} is \(\alpha_\infty(j) := \frac{d}{dt}\operatorname{vol}\{x \in L \otimes \R : Q(x) \le t\}\big|_{t=j}\), the volume being normalized so that \(L\) has covolume \(1\).  In this language Siegel's theorem reads
\begin{equation}\label{eq:siegel-local}
  \frac{1}{\mu_m}\sum_{L \in C_m} \frac{r_L(j)}{g_L} = \alpha_\infty(j) \prod_p \alpha_p(j) \qquad (j \ge 1),
\end{equation}
and the densities depend only on \(L \otimes \Z_p\) and \(L \otimes \R\), which are the same for every \(L \in C_m\).  For every prime \(p\), the reduction \(L/pL\) with the form \(Q \bmod p\) is isometric to the orthogonal sum \(H^k\) of \(k = m/2\) hyperbolic planes, \(Q = x_1y_1 + \dots + x_k y_k\) \cite[Ch.~8]{Ca}.

\begin{problem}\label{pr:alpha-inf}
\(\alpha_\infty(j) = (2\pi)^k j^{k-1}/\Gamma(k)\).
\end{problem}

\begin{problem}\label{pr:count-Fp}
The number of \(x \in \F_p^{2k}\) with \(x_1y_1 + \dots + x_ky_k = j\) is \(p^{2k-1} - p^{k-1}\) for \(j \ne 0\), and \(p^{2k-1} + p^k - p^{k-1}\) for \(j = 0\).
\end{problem}

\begin{problem}\label{pr:alpha-unit}
If \(p \nmid j\), then \(\alpha_p(j) = 1 - p^{-k}\).
\end{problem}

\begin{problem}\label{pr:alpha-p}
If \(p \Vert j\), then \(\alpha_p(j) = (1 - p^{-k})(1 + p^{1-k})\).
\end{problem}

The same argument, splitting off the solutions with \(x \equiv 0 \pmod p\), gives \(\alpha_p(j) = (1 - p^{-k})(1 + p^{1-k}) + p^{2-2k}\alpha_p(j/p^2)\) when \(p^2 \mid j\), the last term coming from \(x = px'\) with \(Q(x') \equiv j/p^2\).  By induction on \(v = v_p(j)\), \(\alpha_p(j) = (1 - p^{-k})\sum_{i=0}^{v} p^{i(1-k)}\).

\begin{problem}\label{pr:siegel-product}
\(\alpha_\infty(j)\prod_p \alpha_p(j) = -\frac{2k}{B_k}\,\sigma_{k-1}(j)\).
\end{problem}

Thus the right side of \eqref{eq:siegel-local} is the \(j\)th coefficient of \(E_k\), and \eqref{eq:siegel-local} is \eqref{eq:siegel} coefficient by coefficient.  With Problems~\ref{pr:E8} and~\ref{pr:rank16} this proves the following.

\begin{theorem}\label{thm:rank8-16}
For \(m = 8\) and \(m = 16\), and every \(j \ge 1\), the mass-weighted mean of \(r_L(j)\) over \(C_m\) is \(\alpha_\infty(j)\prod_p\alpha_p(j)\).
\end{theorem}

Weil reformulated \eqref{eq:siegel} as an identity on adelic groups: the integral of the theta function over \(O(V)(\Q)\backslash O(V)(\mathbb{A})\) taken with the invariant measure of total mass one, is an Eisenstein series, under hypotheses that include the absolute convergence of the Eisenstein series \cite[Theorem~5]{We65}.  For a positive definite space the quotient is compact, and the integral is the average over the genus.  This is the \emph{Siegel--Weil formula in the convergent range}.


\section{The Hermitian class set over \texorpdfstring{\(\Z[\ii]\)}{Z[i]}}\label{sec:hermitian}

Return to the class sets \(\Cl_n\) of Hermitian lattices over \(\Oc = \Z[\ii]\) from Chapter~\ref{dt-chap}.  For \(L \in \Cl_n\) let \(\theta_L(\tau) := \sum_{x \in L} q^{h(x,x)}\), and let \(\theta := \sum_{j \in \Z} q^{j^2}\).  Let \(\chi\) be the non-trivial character modulo \(4\), and \(\Mf_n(4,\chi^n)\) the space of modular forms of weight \(n\) on \(\Gamma_0(4)\) with character \(\chi^n\).  It has dimension \(1 + \lfloor n/2\rfloor\) and basis \(\theta^{2n-4i}F^i\), \(0 \le i \le n/2\), where \(F := \sum_{j \text{ odd}}\sigma_1(j)q^j\) \cite[Ch.~IV, \S1]{Ko}; so an element is determined by its coefficients of \(q^0, \dots, q^{\lfloor n/2\rfloor}\).  Each \(\theta_L\) lies in \(\Mf_n(4,\chi^n)\).

For Hermitian lattices Siegel's theorem is due to Braun \cite{Br}, and it is a case of Weil's theorem \cite[Theorem~5]{We65}; for the unitary formulation see \cite[Theorem~2.3.1]{Li24}.  It reads, for \(n \ge 3\),
\begin{equation}\label{eq:sw-herm}
  \sum_{L \in \Cl_n} \frac{\theta_L}{\lvert \Gamma_L\rvert} = M_n\, E^{(n)}, \qquad E^{(n)} = 1 + \sum_{j \ge 1} \alpha_\infty(j)\prod_p \alpha_p(j)\, q^j ,
\end{equation}
with the densities of \(I_n\) defined as in \S\ref{sec:densities} with \(Q(x) := h(x,x)\), and \(E^{(n)}\) an Eisenstein series in \(\Mf_n(4,\chi^n)\).  The condition \(n \ge 3\) is where Eisenstein series of weight \(n\) converge absolutely.

\begin{problem}\label{pr:converge}
\(\sum_{(c,d) \ne (0,0)} \lvert c\tau + d\rvert^{-s}\) converges if and only if \(s > 2\).
\end{problem}

\begin{problem}\label{pr:theta-L}
\(\theta_{I_n} = \theta^{2n}\), and \(\theta_{E \perp I_1}(\tau) = E_4(2\tau)\,\theta(\tau)^2\), with \(E\) as in Problem~\ref{dt-pr:rank-five}.
\end{problem}

For \(n = 3\) and \(4\), \(\Cl_n = \{I_n\}\) (Problem~\ref{dt-pr:class-one}), so \eqref{eq:sw-herm} says that \(\theta^{2n}\) is itself an Eisenstein series.  This is Jacobi's formula for sums of \(2n\) squares; for \(n = 4\) and odd \(j\) it reads \(r_8(j) = 16\,\sigma_3(j)\).  Here are its local factors.

\begin{problem}\label{pr:r8-odd}
For \(I_4\) and odd \(j\), \(\alpha_\infty(j) = \pi^4j^3/6\) and \(\alpha_\infty(j)\prod_{p \text{ odd}}\alpha_p(j) = 16\,\sigma_3(j)\).
\end{problem}

\begin{problem}\label{pr:mod8}
For odd \(j\), the number of \(x \in (\Z/8\Z)^8\) with \(x_1^2 + \dots + x_8^2 \equiv j \pmod 8\) is \(8^7\).
\end{problem}

Solutions modulo \(8\) with some \(x_i\) odd lift uniformly to \(\Z/2^r\Z\), so Problem~\ref{pr:mod8} gives \(\alpha_2(j) = 1\), and the product of all the local densities is \(16\,\sigma_3(j) = r_8(j)\).

In rank \(5\) the class set has two elements, \(I_5\) and \(E \perp I_1\) (Problem~\ref{dt-pr:rank-five}), and \(\dim \Mf_5(4,\chi) = 3\), one more than the number of Eisenstein series.  Let \(e_4(j) := \sum_{d \mid j,\, d \text{ odd}} \chi(d)\, d^4\) and \(e_4'(j) := \sum_{d \mid j,\, j/d \text{ odd}} \chi(j/d)\, d^4\).  The series \(G_5 := 1 + \frac45\sum_{j \ge 1}(e_4(j) + 16e_4'(j))q^j\) is a combination of the two Eisenstein series of \(\Mf_5(4,\chi)\) \cite[\S4.5]{DS}.

\begin{problem}\label{pr:weights5}
In \(\Cl_5\), the weights \(\lvert\Gamma_L\rvert^{-1}/M_5\) are \(3/5\) for \(I_5\) and \(2/5\) for \(E \perp I_1\).
\end{problem}

\begin{problem}\label{pr:SW5}
\(\frac35\,\theta^{10} + \frac25\,E_4(2\tau)\theta^2 = G_5\).  The coefficients of \(q\) are \(20\), \(4\) and \(68/5\).
\end{problem}

So in rank \(5\) neither theta series is an Eisenstein series, but their mass-weighted average is, as \eqref{eq:sw-herm} predicts.  What is left over is a cusp form.  Let \(f := \frac14\sum_{\alpha \in \Oc}\alpha^4 q^{\lvert\alpha\rvert^2} = q - 4q^2 + 16q^4 - 14q^5 + \cdots\), the theta series of \(\Oc\) with the harmonic polynomial \(\alpha^4\); it is a cusp form in \(\Mf_5(4,\chi)\) \cite[Ch.~10]{Iw}.

\begin{problem}\label{pr:cusp5}
\(\theta^{10} - E_4(2\tau)\theta^2 = 16f\).  Hence \(\theta^{10} = G_5 + \frac{32}{5}f\) and \(E_4(2\tau)\theta^2 = G_5 - \frac{48}{5}f\).
\end{problem}

The first identity of the last line is Liouville's formula \(r_{10}(j) = \frac45(e_4(j) + 16e_4'(j)) + \frac85\sum_{\lvert\alpha\rvert^2 = j}\alpha^4\).  The cusp parts \(\frac{32}{5}f\) and \(-\frac{48}{5}f\) cancel in the average with weights \(3/5\) and \(2/5\), and only there.


\section{Gram rows and nonvanishing}\label{sec:gram}

The theta lifts of the argument go from a class set \(\Cl_n\) to a much larger class set \(\Cl_{n'}\), of which only the mass is known (Chapter~\ref{dt-chap}, \S\ref{dt-sec:mass}).  Their Gram matrices are sums over \(\Cl_{n'}\).  We look at the smallest case, \(n = 1\).

For \(a \ge 0\) and \(L' \in \Cl_{n'}\) put
\[
  v_a(L') := \sum_{x \in L'} x^a\, e^{-\pi h(x,x)} \in \Sym^a(L'_\C),
\]
with \(x^a\) as in \S\ref{dt-sec:forms}.  Up to the factor \(\lvert\Gamma_{I_1}\rvert = 4\), this is the theta lift of the next chapter, \emph{Definite theta lifts and their norms}, applied to the form on \(\Cl_1 = \{I_1\}\) whose value is the basis vector of \(\Sym^a(\C)\).  For \(x, y \in L'\) let \(\Gram(x,y)\) be the matrix with entries \(h(x,x)\), \(h(x,y)\), \(h(y,x)\), \(h(y,y)\), and for a \(2\times2\) Hermitian matrix \(T\) over \(\Oc\) put
\[
  R_{n'}(T) := \sum_{L' \in \Cl_{n'}} \frac{1}{\lvert\Gamma_{L'}\rvert}\, \#\{(x,y) \in L'^2 : \Gram(x,y) = T\} .
\]

\begin{problem}\label{pr:va-form}
\(v_a\) is a form of weight \((a)\) on \(\Cl_{n'}\), and \(v_a = 0\) unless \(4 \mid a\).
\end{problem}

\begin{problem}\label{pr:gram}
\(\lVert v_a\rVert^2 = \sum_T R_{n'}(T)\, T_{12}^a\, e^{-\pi \operatorname{tr} T}\), the sum running over positive semi-definite \(T\).
\end{problem}

So the norm of the lift, and in the same way every entry of its Gram matrix, depends on \(\Cl_{n'}\) only through the averaged representation numbers \(R_{n'}(T)\).  The Siegel--Weil formula of degree \(2\) evaluates them: for \(n' > 4\) and \(T\) positive definite, \(R_{n'}(T) = M_{n'}\,\alpha_\infty(T)\prod_p\alpha_p(T)\), where \(\alpha_v(T)\) are the local densities of representations of \(T\) by \(I_{n'}\), defined as in \S\ref{sec:densities} with pairs of vectors in place of vectors \cite[Theorem~5]{We65}.  Here \(n' > 4\) is the convergent range for pairs, as \(n \ge 3\) was for single vectors.  In particular \(R_{n'}(T) > 0\) as soon as \(T\) is represented by \(I_{n'}\) over every \(\Z_p\): a statement about the average that need not hold for any single lattice.

\begin{problem}\label{pr:identity2}
\(I_5\) contains \(320\) pairs \((x,y)\) with \(\Gram(x,y) = 1\), and \(E \perp I_1\) contains none.  Hence \(R_5(1) = 1/384\).
\end{problem}

For \(a = 0\) one can take the lift one step further and pair it with the constant form.

\begin{problem}\label{pr:period}
For \(n' \ge 3\), \(\langle v_0, 1\rangle = M_{n'}\,E^{(n')}(\ii/2)\).  For \(n' = 5\) this is \(G_5(\ii/2)/73728\), and \(G_5(\ii/2) = 2.06195\ldots\).
\end{problem}

The general statement is the Rallis inner product formula, which we state without proof.  Let \(\varphi\) be a form on \(\Cl_n\) of weight \(\lambda\) that is an eigenvector of the Hecke operators at almost all primes, and let \(\Theta\varphi\) be its theta lift to \(\Cl_{n'}\) with \(n' > 4n\).  Then \(\lVert\Theta\varphi\rVert^2/\lVert\varphi\rVert^2\) is the product of finitely many local factors and of a partial \(L\)-function of \(\varphi\) at a point of its half-plane of absolute convergence.  The proof unfolds \(\lVert\Theta\varphi\rVert^2\) into an integral of \(\varphi \otimes \bar\varphi\) against the genus average of the doubled theta series, the analogue of Problem~\ref{pr:gram}, and evaluates that average by the Siegel--Weil formula \cite{Ra}; for unitary groups in the stable range see \cite{JSLi}.  Since an Euler product does not vanish where it converges absolutely, the \(L\)-value is nonzero, and the lift is nonzero once the local factors are.

In Adriazola's reconstruction of the argument the convergent-range Siegel--Weil formula is the one external analytic theorem behind the Gram-fibre identities \cite[\S14.4]{Adr}; the near orthogonality that these identities feed into rests there also on counting and pre-trace statements, \cite[Hyps.~14.61 and 14.71--14.77]{Adr}, assumed there (see Chapter~\ref{ps-chap}, and the ledger in Section~\ref{as-sec:map} of Chapter~\ref{as-chap}).  Further, the norms of the lifts enter through the Rallis inner product formula in the convergent range \cite[Prop.~15.52(i)]{Adr}; \cite[\S15.13]{Adr} specifies the forms on a selected realization, and the formula for the lifts of the chain is assumed there.  In both places the class sets are far too large to handle one lattice at a time, and only the average, through its local densities, is used.


\section*{Hints}

\subsection*{\S\ref{sec:even}}
\begin{description}
\item[\ref{pr:E8}] \(\dim\Mf_4 = 1\).
\item[\ref{pr:rank16}] \(\dim \Mf_8 = 1\), and \(\theta_{E_8\oplus E_8} = \theta_{E_8}^2\).  For the roots of \(D_{16}^+\), compute \(x\cdot x\) on the non-trivial coset.
\item[\ref{pr:not-iso}] Is each lattice spanned by its roots?
\item[\ref{pr:leech}] Write \(\theta_\Lambda = E_{12} + c\Delta\) and compare coefficients of \(q\); then use \(\tau(2) = -24\).
\item[\ref{pr:mass}] \(B_2, B_4, \dots, B_{16}\); for \(m = 16\) one needs \(B_8 = -1/30\) and \(B_{14} = 7/6\).
\item[\ref{pr:niemeier}] Subtract \(\mu_{24}E_{12}\) from both sides of \eqref{eq:siegel}, and compare coefficients of \(q\).
\end{description}

\subsection*{\S\ref{sec:densities}}
\begin{description}
\item[\ref{pr:alpha-inf}] Since \(L\) is unimodular, the volume is that of a Euclidean ball of radius \(\sqrt{2t}\) in dimension \(2k\).
\item[\ref{pr:count-Fp}] Induction on \(k\): count first the solutions with \(x_k = 0\), then those with \(x_k \ne 0\).
\item[\ref{pr:alpha-unit}] Every solution modulo \(p\) of \(Q(x) \equiv j\) has \(x \not\equiv 0\), and the gradient of \(Q\) is nonzero there; apply Hensel's lemma.
\item[\ref{pr:alpha-p}] Separate the solutions with \(x \not\equiv 0 \pmod p\), which lift as in Problem~\ref{pr:alpha-unit}, from those with \(x \equiv 0\), for which \(Q(x) \equiv 0 \pmod{p^2}\).
\item[\ref{pr:siegel-product}] The Euler product of \(\zeta(k)\), Euler's formula for \(\zeta(k)\) in terms of \(B_k\), and \(\sigma_{k-1}(j) = j^{k-1}\sigma_{1-k}(j)\).
\end{description}

\subsection*{\S\ref{sec:hermitian}}
\begin{description}
\item[\ref{pr:converge}] Count the pairs with \(N \le \lvert c\tau + d\rvert < N+1\).
\item[\ref{pr:theta-L}] \(h(x,x)\) is additive on orthogonal sums, and \(h(x,x) = x\cdot x\) is even on \(E\).
\item[\ref{pr:r8-odd}] As in Problem~\ref{pr:alpha-inf}, but \(\{Q \le t\}\) is now the ball of radius \(\sqrt t\).  The counting argument of Problems~\ref{pr:alpha-unit}--\ref{pr:alpha-p} and the remark after them uses only that \(Q \bmod p\) is isometric to \(H^4\), which holds for \(x_1^2 + \dots + x_8^2\) and odd \(p\).  Then \(\zeta(4) = \pi^4/90\).
\item[\ref{pr:mod8}] Write the count as \(\frac18\sum_{t \bmod 8} e(-tj/8)\,g(t)^8\) with the Gauss sums \(g(t) := \sum_{x \bmod 8} e(tx^2/8)\).
\item[\ref{pr:weights5}] Problem~\ref{dt-pr:aut} and \(\lvert\Gamma_E\rvert = 46080\).
\item[\ref{pr:SW5}] Both sides lie in \(\Mf_5(4,\chi)\); compare the coefficients of \(q^0\), \(q^1\), \(q^2\).
\item[\ref{pr:cusp5}] As in Problem~\ref{pr:SW5}; then combine with it.
\end{description}

\subsection*{\S\ref{sec:gram}}
\begin{description}
\item[\ref{pr:va-form}] Substitute \(x \mapsto \gamma x\) for \(\gamma \in \Gamma_{L'}\), and \(x \mapsto \ii x\).
\item[\ref{pr:gram}] Problem~\ref{dt-pr:sym-norm}: \(\langle x^a, y^a\rangle = h(x,y)^a\).  Then group the pairs by their Gram matrix.
\item[\ref{pr:identity2}] Problem~\ref{dt-pr:rank-five}: the vectors of \(E \perp I_1\) with \(h(x,x) = 1\) are \(ue\), \(u \in \mu_4\).
\item[\ref{pr:period}] What is \(v_0(L')\) as a value of \(\theta_{L'}\)?
\end{description}

\begin{thebibliography}{JSL}
\bibitem[Adr]{Adr} J.~Adriazola, \emph{The Alp\"oge proof of the Riemann hypothesis}, manuscript, revised 25 November 2026, \href{https://github.com/jimjamzola/RH-Alpoge-Reconstruction/tree/9eefc08b0819728d4dd2ef0568b49a61251af111}{commit \texttt{9eefc08}}.
\bibitem[Br]{Br} H.~Braun, \emph{Zur Theorie der hermitischen Formen}, Abh.\ Math.\ Sem.\ Hansischen Univ.\ \textbf{14} (1941), 61--150.
\bibitem[Ca]{Ca} J.~W.~S.~Cassels, \emph{Rational Quadratic Forms}, Academic Press, 1978.
\bibitem[CoS]{CoS} J.~H.~Conway and N.~J.~A.~Sloane, \emph{Sphere Packings, Lattices and Groups}, 3rd ed., Grundlehren Math.\ Wiss.\ 290, Springer, 1999.
\bibitem[DS]{DS} F.~Diamond and J.~Shurman, \emph{A First Course in Modular Forms}, Grad.\ Texts in Math.\ 228, Springer, 2005.
\bibitem[Iw]{Iw} H.~Iwaniec, \emph{Topics in Classical Automorphic Forms}, Grad.\ Stud.\ Math.\ 17, Amer.\ Math.\ Soc., 1997.
\bibitem[Ko]{Ko} N.~Koblitz, \emph{Introduction to Elliptic Curves and Modular Forms}, 2nd ed., Grad.\ Texts in Math.\ 97, Springer, 1993.
\bibitem[JSL]{JSLi} J.-S.~Li, \emph{Non-vanishing theorems for the cohomology of certain arithmetic quotients}, J.\ Reine Angew.\ Math.\ \textbf{428} (1992), 177--217.
\bibitem[Li]{Li24} C.~Li, \emph{Geometric and arithmetic theta correspondences}, arXiv:2402.12159, 2024.
\bibitem[Mi]{Mi} J.~Milnor, \emph{Eigenvalues of the Laplace operator on certain manifolds}, Proc.\ Nat.\ Acad.\ Sci.\ USA \textbf{51} (1964), 542.
\bibitem[Ra]{Ra} S.~Rallis, \emph{Injectivity properties of liftings associated to Weil representations}, Compositio Math.\ \textbf{52} (1984), 139--169.
\bibitem[SeA]{SeA} J.-P.~Serre, \emph{A Course in Arithmetic}, Grad.\ Texts in Math.\ 7, Springer, 1973.
\bibitem[Si]{Si} C.~L.~Siegel, \emph{\"Uber die analytische Theorie der quadratischen Formen}, Ann.\ of Math.\ \textbf{36} (1935), 527--606.
\bibitem[We65]{We65} A.~Weil, \emph{Sur la formule de Siegel dans la th\'eorie des groupes classiques}, Acta Math.\ \textbf{113} (1965), 1--87.
\end{thebibliography}

\end{document}
